\chapter{基础计算：Outs、胜率与赔率}

\section{本章导入}

前面几章中，我们多次提到“听牌是否值得继续”“面对下注是否有足够胜率”“下注尺度会影响对手继续范围”“河牌跟注必须满足数学条件”。这些内容背后，都离不开一个最基础的问题：

\begin{center}
\textbf{继续这手牌，值不值当前这个价格？}
\end{center}

很多新手在牌桌上判断是否跟注时，依靠的是感觉。他们会想：“我还有同花机会”“我还有顺子机会”“他可能在 bluff”“都打到这里了，再看一张”。这些想法都有一个共同问题：它们只说明你还有机会，却没有说明这个机会是否值得你付出当前成本。

德州扑克不是只看有没有机会，而是看机会和价格是否匹配。
同花听牌当然有机会完成，但如果对手下注太大，你付出的价格可能不划算。
顶对当然有机会领先，但如果对手河牌大注范围几乎全是强牌，你的胜率可能不够。
一手牌当然有可能赢，但如果你长期赢不到足够比例，跟注就是亏损。

因此，在进入下一章“弃牌与跟注”之前，我们需要先补上一套基础计算工具。本章要讲三个最重要的入门概念：

\begin{itemize}
\item Outs：哪些未发出的牌能帮助你改善；
\item 胜率：你的牌最终赢下底池的大致概率；
\item 赔率：你为了继续游戏需要付出的价格是否划算。
\end{itemize}

本章不是数学课，也不要求你在牌桌上进行复杂精确计算。你只需要掌握几条简单但非常实用的规则：如何数 outs，如何用“2/4 法则”快速估算胜率，如何计算底池赔率，如何比较“实际胜率”和“所需胜率”。

本章的核心输出是“继续决策四步法”：

\begin{center}
\textbf{先数 outs。\\
再估胜率。\\
计算跟注所需胜率。\\
比较实际胜率与所需胜率。}
\end{center}

学完本章之后，你会明白一个最重要的原则：

\begin{center}
\textbf{能不能跟，不看“有没有希望”，而看“这个希望值不值这个价格”。}
\end{center}

\section{核心概念}

\subsection{什么是 Outs}

Outs 可以理解为：牌堆中还没有发出的牌里，哪些牌能让你的手牌改善为大概率领先的牌。中文里可以把它理解为“补牌”或“有效补牌”。

例如你手里是：

\begin{center}
\texttt{As 7s}
\end{center}

翻牌是：

\begin{center}
\texttt{Ks 9s 2d}
\end{center}

你现在有同花听牌。只要转牌或河牌再来一张黑桃，你就能组成同花。

一副牌中每种花色有 13 张牌。你已经看到 4 张黑桃：

\begin{center}
\texttt{As、7s、Ks、9s}
\end{center}

所以还剩：

\[
13 - 4 = 9
\]

张黑桃没有出现。这 9 张黑桃，就是你的 outs。

也就是说：

\begin{center}
\textbf{同花听牌通常有 9 个 outs。}
\end{center}

再看一个顺子听牌的例子。你手里是：

\begin{center}
\texttt{9c 8c}
\end{center}

翻牌是：

\begin{center}
\texttt{7d 6s Kh}
\end{center}

现在只要转牌或河牌出现 \texttt{5} 或 \texttt{T}，你就能组成顺子。一副牌中有 4 张 5 和 4 张 T，因此你一共有：

\[
4 + 4 = 8
\]

个 outs。

这就是两头顺听牌，通常有 8 个 outs。

Outs 的意义在于，它把“我还有机会”变成了更具体的数字。你不再只是模糊地说“我有听牌”，而是能够说清楚“我有 9 个 outs”或“我有 8 个 outs”。

\subsection{常见 Outs 类型}

新手可以先记住几类最常见的 outs。

\begin{center}
\begin{tabular}{lll}
\toprule
\textbf{听牌类型} & \textbf{典型情况} & \textbf{常见 outs} \\
\midrule
同花听牌 & 四张同花，差一张成同花 & 9 \\
两头顺听牌 & 两端都能成顺 & 8 \\
卡顺听牌 & 只有中间一张能成顺 & 4 \\
两张高牌成一对 & 两张高牌都可能配对 & 约 6 \\
一对变三条 & 口袋对子或一对想中三条 & 2 \\
同花听牌加两头顺 & 组合听牌 & 约 15 \\
\bottomrule
\end{tabular}
\end{center}

例如：

\begin{itemize}
\item 同花听牌：通常 9 outs；
\item 两头顺听牌：通常 8 outs；
\item 卡顺听牌：通常 4 outs；
\item 两张高牌想击中一对：通常最多 6 outs；
\item 一对想变三条：通常 2 outs；
\item 同花听牌加两头顺：可能有 15 outs 左右。
\end{itemize}

这里要注意，“常见 outs”只是入门估算。实战中，你还要判断这些 outs 是否真的有效。有些牌虽然能让你成牌，但可能也会让对手成更强牌。这种 outs 不能全部算满。

\subsection{干净 Outs 与不干净 Outs}

不是所有 outs 都是干净的。

干净 outs 指的是：如果这张牌出现，你大概率能成为最好牌。
不干净 outs 指的是：这张牌虽然改善了你，但也可能让你输给更强牌。

例如你手里是：

\begin{center}
\texttt{As 7s}
\end{center}

翻牌是：

\begin{center}
\texttt{Ks 9s 2d}
\end{center}

你追的是 A 高同花。如果后面来黑桃，你大概率能组成很强的同花。因为你手中有 \texttt{As}，你的同花是坚果同花。这类 outs 比较干净。

但如果你手里是：

\begin{center}
\texttt{8s 7s}
\end{center}

翻牌是：

\begin{center}
\texttt{Ks Qs 2d}
\end{center}

你同样有同花听牌，但即使后面来黑桃，你也只是 8 高同花。如果对手持有 \texttt{AsJs}、\texttt{AsTs}、\texttt{JsTs} 等更高同花，你完成同花后仍然可能输掉大底池。

这就是不干净 outs。

再看顺子听牌。你如果追的是低端顺子，而牌面可能同时让对手组成更高顺子，那么部分 outs 也要打折。比如你成了小顺子，但对手可能成更大顺子，这就是反向风险。

因此，新手不能机械地认为“我有 9 个 outs，所以一定可以跟”。你还要问：

\begin{itemize}
\item 我完成后是否接近坚果？
\item 我的 outs 是否可能让对手也完成更强牌？
\item 我的同花是否可能被更高同花压制？
\item 我的顺子是否可能输给更大顺子？
\item 我的高牌配对后是否可能输给更好踢脚？
\end{itemize}

如果 outs 不干净，就要适当打折。你可以不必精确计算成 7.5 个 outs 或 6.2 个 outs，但必须知道：低质量听牌的真实价值低于表面数字。

\subsection{什么是胜率}

胜率指的是：从当前局面到摊牌，你最终赢下底池的概率。

在听牌场景中，我们通常先用 outs 来估算胜率。例如你有 9 个 outs 的同花听牌，就可以估算你在转牌或河牌完成同花的大致概率。

新手不需要一开始掌握复杂概率公式，只需要先记住一个非常实用的近似方法：

\begin{center}
\textbf{翻牌到河牌：outs $\times 4 \approx$ 成牌概率。\\
转牌到河牌：outs $\times 2 \approx$ 成牌概率。}
\end{center}

这就是常说的“2/4 法则”。

例如你在翻牌有同花听牌，通常是 9 个 outs。翻牌之后还有转牌和河牌两张牌可以发，因此：

\[
9 \times 4 = 36\%
\]

也就是说，翻牌同花听牌到河牌完成同花的概率大约是 36\%。实际精确值约为 35\%，这个估算已经足够实战使用。

如果到了转牌，你仍然只有同花听牌，只剩河牌一张牌：

\[
9 \times 2 = 18\%
\]

实际精确值约为 19.6\%。同样，这个估算已经足够帮助你做初步决策。

\subsection{2/4 法则}

2/4 法则是新手最实用的胜率估算方法。

\begin{center}
\begin{tabular}{lll}
\toprule
\textbf{当前街道} & \textbf{还剩公共牌} & \textbf{估算方法} \\
\midrule
翻牌圈 & 转牌和河牌两张 & outs $\times 4$ \\
转牌圈 & 河牌一张 & outs $\times 2$ \\
\bottomrule
\end{tabular}
\end{center}

例如：

\begin{itemize}
\item 同花听牌 9 outs：翻牌约 36\%，转牌约 18\%；
\item 两头顺听牌 8 outs：翻牌约 32\%，转牌约 16\%；
\item 卡顺听牌 4 outs：翻牌约 16\%，转牌约 8\%；
\item 一对变三条 2 outs：翻牌约 8\%，转牌约 4\%；
\item 强组合听牌 15 outs：翻牌约 60\%，转牌约 30\%。
\end{itemize}

这条规则有两个重要作用。

第一，它让你快速知道听牌大概有多少胜率。你不再只是说“我有同花机会”，而是能说“我大约有 36\% 的完成概率”。

第二，它让你能和底池赔率进行比较。你可以判断当前跟注价格是否划算。

需要注意的是，2/4 法则估算的是“改善成目标牌”的概率，而不一定等于“最终赢牌概率”。因为你完成听牌后仍可能输给更强牌，你没完成听牌时也可能偶尔靠高牌或对子赢。因此它是实用估算，而不是绝对答案。

\subsection{什么是底池赔率}

底池赔率回答的问题是：

\begin{center}
\textbf{面对当前下注，我跟注需要赢多少比例，才不亏？}
\end{center}

计算公式是：

\[
\text{所需胜率} =
\frac{\text{跟注金额}}{\text{跟注后总底池}}
\]

这里要特别注意，是“跟注后总底池”，不是当前底池。

例如，当前底池为 100，对手下注 50。你需要跟注 50。若你跟注，底池中会有：

\[
100 + 50 + 50 = 200
\]

其中包括原底池 100，对手下注 50，以及你的跟注 50。

因此，你需要的胜率是：

\[
\frac{50}{200} = 25\%
\]

也就是说，只要你长期能赢超过 25\%，这个跟注才可能盈利。如果你的实际胜率低于 25\%，跟注就是亏损。

再看一个例子。当前底池为 100，对手下注 100。你要跟注 100。跟注后总底池为：

\[
100 + 100 + 100 = 300
\]

所需胜率为：

\[
\frac{100}{300} = 33.3\%
\]

也就是说，面对一个底池大小下注，你大约需要 33.3\% 的胜率。

底池赔率的核心不是让你背公式，而是让你建立价格意识：

\begin{center}
\textbf{小注需要的胜率低，大注需要的胜率高。}
\end{center}

\subsection{常见下注尺度对应的所需胜率}

新手可以先记住几个常见下注尺度对应的所需胜率。

\begin{center}
\begin{tabular}{llll}
\toprule
\textbf{对手下注} & \textbf{原底池} & \textbf{跟注后总底池} & \textbf{所需胜率} \\
\midrule
1/3 底池 & 100 & 166.7 & 约 20\% \\
1/2 底池 & 100 & 200 & 25\% \\
2/3 底池 & 100 & 233.3 & 约 28.6\% \\
1 倍底池 & 100 & 300 & 33.3\% \\
1.5 倍底池 & 100 & 400 & 37.5\% \\
\bottomrule
\end{tabular}
\end{center}

这张表非常重要。它会直接影响你面对下注时的防守范围。

面对 1/3 底池小注，你只需要约 20\% 胜率，因此可以用较宽范围继续。
面对 1/2 底池下注，你需要 25\% 胜率。
面对底池大小下注，你需要约 33.3\% 胜率。
面对超池下注，你需要的胜率更高，必须明显收紧跟注范围。

很多新手的问题是，不管对手下注多大，他们都用差不多的牌去跟。这是严重错误。下注尺度越大，你继续所需的胜率越高，对手范围通常也越强，你必须更加谨慎。

\subsection{需要胜率与实际胜率}

跟注决策中，有两个概念必须区分：需要胜率和实际胜率。

需要胜率，是由底池赔率决定的。
实际胜率，是你的手牌面对对手范围时，实际能赢的概率。

比如当前底池 100，对手下注 50。你跟注需要 25\% 胜率。这是需要胜率。

但你到底有没有 25\% 胜率，要看你的牌面对手下注范围的表现。如果你有同花听牌，可能有大约 18\% 到 36\% 的改善概率，具体看是在翻牌还是转牌。如果你在河牌拿着中等对子面对对手大注，你可能只能击败诈唬，实际胜率取决于对手诈唬比例。

很多新手只看需要胜率，却不判断实际胜率。例如他们会说：“这个下注不大，我赔率不错，可以跟。”但如果对手几乎全是强价值牌，你实际胜率远远不够，再好的价格也不能乱跟。

也有相反情况。对手下注较大，你所需胜率较高，但如果你有强组合听牌，实际胜率也可能足够，甚至可以主动加注半诈唬。

所以，正确的比较是：

\begin{center}
\textbf{实际胜率 $\geq$ 所需胜率：可以考虑继续。\\
实际胜率 $<$ 所需胜率：通常应该弃牌。}
\end{center}

这里说“可以考虑继续”，不是说一定要跟。因为实战还要考虑位置、对手类型、隐含赔率、反向隐含赔率和后续行动。

\subsection{隐含赔率}

隐含赔率指的是：你现在跟注虽然从直接底池赔率看不一定完全划算，但如果后续成牌，可能还能从对手身上赢到更多筹码。

最典型例子是小对子追暗三条。

你翻前用 \texttt{55} 跟注对手加注。你现在的目标不是靠 5 高对子摊牌赢，而是希望翻牌击中暗三条。如果你中了暗三条，而对手持有 \texttt{AA}、\texttt{KK} 或强顶对，他可能会支付很多筹码。你当前跟注的价值，部分来自未来可能赢到的大底池，这就是隐含赔率。

同花听牌和顺子听牌也可能有隐含赔率。你现在跟注追听牌，如果后面完成强牌，对手愿意继续支付，你的跟注价值就会上升。

但隐含赔率成立需要条件：

\begin{itemize}
\item 有效筹码足够深；
\item 你完成后牌力足够强；
\item 对手会在你成牌后继续支付；
\item 你的成牌不明显，不容易让对手立刻停止下注；
\item 你有位置，方便控制后续行动。
\end{itemize}

如果筹码很浅，即使你成牌也赢不到更多，隐含赔率就低。
如果你追的是明显同花，而第三张同花一出现对手立刻不再支付，隐含赔率也低。
如果你追的是低同花，完成后可能输给更高同花，隐含赔率看似高，实则包含反向风险。

因此，隐含赔率不是“我中了可以赢很多”的幻想，而是要判断对手是否真的会在你成牌后付钱。

\subsection{反向隐含赔率}

反向隐含赔率是新手非常容易忽略的概念。它指的是：你成牌后看似变强了，但实际上可能输给更强牌，结果反而输掉更大的底池。

典型例子包括：

\begin{itemize}
\item 弱 A 击中 A，但输给更强 A；
\item 低同花完成同花，但输给更高同花；
\item 小顺子完成顺子，但输给更大顺子；
\item 非坚果听牌完成后，被对手更强听牌压制；
\item 两张高牌配对后，输给更好踢脚。
\end{itemize}

例如你持有：

\begin{center}
\texttt{Ac 6d}
\end{center}

翻牌为：

\begin{center}
\texttt{Ah Js 8c}
\end{center}

你击中了顶对，但踢脚很弱。如果对手范围中有很多 \texttt{AK}、\texttt{AQ}、\texttt{AJ}、\texttt{AT}，你虽然“中了 A”，却很容易输给更强 A。此时你的牌有很强的反向隐含赔率。

再如你持有：

\begin{center}
\texttt{8s 7s}
\end{center}

牌面为：

\begin{center}
\texttt{Ks Qs 2d}
\end{center}

你追同花，但即使完成，也可能输给更高同花。你看似有 9 个 outs，但这些 outs 并不完全干净。

反向隐含赔率提醒我们：

\begin{center}
\textbf{不是所有“成牌”都值得追，不是所有“中了”都舒服。}
\end{center}

很多边缘牌和弱听牌的问题，不是它们完全没有机会，而是它们在成牌后仍然可能输大底池。

\subsection{把 Outs、胜率和赔率连起来}

真正实战中，outs、胜率和赔率不是孤立使用的。它们应该连成一个完整决策流程。

\begin{enumerate}
\item 先数 outs：哪些牌能帮助我改善？
\item 再估胜率：用 2/4 法则估算大致完成概率；
\item 计算所需胜率：根据底池赔率判断我需要赢多少；
\item 比较两者：实际胜率是否足够？
\item 调整判断：outs 是否干净？是否有隐含赔率？是否有反向隐含赔率？
\end{enumerate}

例如，你翻牌有同花听牌，9 outs。用 2/4 法则估算，翻牌到河牌约有 36\% 完成概率。如果对手下注给你的价格只需要 25\% 胜率，那么从直接赔率看，可以考虑继续。

但如果这是低同花听牌，且对手范围中可能有更高同花，你要给 outs 打折。
如果对手很强，成牌后不一定支付，你的隐含赔率下降。
如果你没有位置，转牌还可能面对更大下注，你也要谨慎。

所以，入门计算不是让你机械跟注，而是让你有一个理性起点。你先用数字判断价格是否大致合理，再结合牌面、范围和对手类型做最终决策。

\section{新手常见误区}

\subsection{误区一：只要有 Outs 就继续}

很多新手看到自己有听牌，就觉得必须跟注。他们会说：“我还有 9 张牌能中同花。”“我还有 8 张牌能成顺。”但他们没有问：对手下注给我的价格是否划算？

有 outs 只是说明你有机会，并不说明这个机会值得购买。面对小注，同花听牌可能轻松继续；面对转牌大注，同样的同花听牌可能就不够价格。

正确思路是：先数 outs，再看价格。没有价格意识的追牌，是长期亏损的重要来源。

\subsection{误区二：把所有 Outs 都算满}

新手常常机械地数 outs，却不判断 outs 是否干净。

低同花听牌不应该和坚果同花听牌等价。
弱 A 击中 A 不应该和强 A 等价。
小顺子听牌不应该和坚果顺子听牌等价。

如果你的 outs 可能让你成第二好牌，就要打折。很多时候，问题不是你中不了，而是你中了以后还会输更多。

\subsection{误区三：翻牌胜率和转牌胜率混淆}

翻牌听牌和转牌听牌差别很大。

翻牌同花听牌还有两张牌，约 36\%。
转牌同花听牌只剩一张牌，约 18\%。

很多新手翻牌跟注合理，但转牌面对大注仍然机械继续，就是因为他们没有意识到胜率已经明显下降。到了转牌，听牌继续条件更加苛刻。

\subsection{误区四：只看赔率，不看对手范围}

有些玩家学会底池赔率后，又走向另一个误区：只要价格看起来不错，就跟注。

但底池赔率只告诉你“需要多少胜率”，不告诉你“实际有没有这么多胜率”。实际胜率必须结合对手范围。

例如河牌面对小注，你可能只需要较低胜率。但如果对手是极度被动玩家，河牌下注几乎全是价值牌，你的中等牌力仍然可能不能跟。

价格重要，范围同样重要。

\subsection{误区五：幻想隐含赔率}

很多新手会说：“我中了就能赢很多。”但这常常只是幻想。

如果对手看到同花完成就不再支付，你的同花听牌没有你想象中那么高隐含赔率。
如果有效筹码很浅，你即使成牌也赢不到更多。
如果你追的是低同花，完成后可能输给更高同花，那么所谓隐含赔率可能变成反向隐含赔率。

隐含赔率必须建立在真实支付可能性上，而不是想象中的大底池。

\subsection{误区六：忽略位置和后续行动}

听牌跟注不仅要看当前赔率，还要看后续行动。

你在有位置时跟注，后续可以根据对手行动决定是否继续。
你在无位置时跟注，转牌可能再次面对下注，无法轻松实现权益。

尤其在翻牌面对下注时，如果你跟注后转牌还可能被大额施压，那么你不能只看当前一街价格。你要考虑整手牌路径。

\section{实战思考框架}

本章的核心框架是“继续决策四步法”：

\begin{center}
\textbf{先数 outs。\\
再估胜率。\\
计算跟注所需胜率。\\
比较实际胜率与所需胜率。}
\end{center}

\subsection{第一步：先数 Outs}

面对下注时，如果你还没有成强牌，而是依靠后续改善继续，第一步就是数 outs。

你要问：

\begin{itemize}
\item 哪些牌能让我成同花？
\item 哪些牌能让我成顺子？
\item 哪些牌能让我从一对变三条？
\item 哪些高牌配对后可能让我领先？
\item 这些 outs 是否真的干净？
\end{itemize}

不要只说“我有听牌”。要尽量具体：我是 9 outs 同花听牌，还是 4 outs 卡顺？我是坚果同花听牌，还是低同花听牌？这些差别会直接影响后续决策。

\subsection{第二步：再估胜率}

数出 outs 后，用 2/4 法则估算胜率。

\begin{itemize}
\item 翻牌到河牌：outs $\times 4$；
\item 转牌到河牌：outs $\times 2$。
\end{itemize}

例如：

\begin{itemize}
\item 翻牌 9 outs：约 36\%；
\item 转牌 9 outs：约 18\%；
\item 翻牌 8 outs：约 32\%；
\item 转牌 8 outs：约 16\%；
\item 翻牌 4 outs：约 16\%；
\item 转牌 4 outs：约 8\%。
\end{itemize}

这里的胜率是估算值，不是绝对精确数字。它的作用是帮助你快速判断“我大概有多少机会”。

\subsection{第三步：计算跟注所需胜率}

然后，根据底池赔率计算你需要多少胜率。

公式是：

\[
\text{所需胜率} =
\frac{\text{跟注金额}}{\text{跟注后总底池}}
\]

如果不想每次精确计算，可以先记住常见尺度：

\begin{center}
\begin{tabular}{ll}
\toprule
\textbf{下注尺度} & \textbf{所需胜率} \\
\midrule
1/3 底池 & 约 20\% \\
1/2 底池 & 25\% \\
2/3 底池 & 约 28.6\% \\
1 倍底池 & 33.3\% \\
1.5 倍底池 & 37.5\% \\
\bottomrule
\end{tabular}
\end{center}

下注越大，你需要的胜率越高。面对大注，用弱听牌继续通常会变得很困难。

\subsection{第四步：比较实际胜率与所需胜率}

最后，把估算胜率和所需胜率进行比较。

如果你的实际胜率高于所需胜率，可以考虑跟注。
如果你的实际胜率低于所需胜率，通常应该弃牌。

例如你翻牌有同花听牌，约 36\% 完成概率。对手下注给你的价格只需要 25\% 胜率，从直接赔率看可以跟注。

但如果你转牌仍然只有同花听牌，约 18\% 完成概率，对手下注底池大小，你需要约 33.3\% 胜率。此时如果没有很强隐含赔率或额外胜率，跟注通常不划算。

\subsection{第五步：用质量因素修正判断}

前四步给出的是基础数学判断。实战中还要修正。

你需要继续问：

\begin{itemize}
\item 我的 outs 是否干净？
\item 我完成后是否接近坚果？
\item 对手是否会在我成牌后继续支付？
\item 我是否有位置？
\item 有效筹码是否足够深？
\item 如果我没中，后续是否还能 bluff？
\item 如果我中了，是否可能输给更强牌？
\end{itemize}

这些因素不会推翻基础计算，但会影响最终决策。强听牌、有位置、深筹码、对手愿意支付，会让继续更有吸引力。弱听牌、无位置、浅筹码、反向隐含赔率高，则会让继续变差。

\section{牌例拆解}

\subsection{牌例一：翻牌同花听牌面对半池下注}

\subsubsection*{牌局背景}

有效筹码 100BB。你在按钮位拿到：

\begin{center}
\texttt{As 7s}
\end{center}

你 open raise 到 2.5BB，大盲跟注。

翻牌为：

\begin{center}
\texttt{Ks 9s 2d}
\end{center}

大盲下注 1/2 底池，行动轮到你。

\subsubsection*{新手常见想法}

新手可能会想：

\begin{itemize}
\item “我有同花听牌，应该跟。”
\item “如果来黑桃，我就是同花。”
\item “A 高同花应该很强。”
\end{itemize}

这些想法方向不错，但还需要用计算确认价格。

\subsubsection*{理性分析}

第一，数 outs。你有 A 高同花听牌。一共有 13 张黑桃，你已经看到 4 张黑桃，因此还有 9 张黑桃可以帮助你成同花。

第二，估胜率。翻牌到河牌还有两张牌，用 2/4 法则：

\[
9 \times 4 = 36\%
\]

你完成同花的概率大约是 36\%。

第三，计算所需胜率。对手下注 1/2 底池。面对半池下注，你大约需要 25\% 胜率。

第四，比较。你的估算完成概率约 36\%，高于所需胜率 25\%。而且你是 A 高同花听牌，outs 较干净。

\subsubsection*{推荐行动}

可以跟注，也可以根据对手类型和弃牌权益考虑加注半诈唬。

如果对手容易弃牌，加注有吸引力；如果对手不爱弃牌，跟注实现权益也合理。

\subsubsection*{本手牌教训}

翻牌坚果同花听牌面对半池下注，通常有足够直接赔率继续。关键不是“有听牌就跟”，而是胜率确实高于所需胜率。

\subsection{牌例二：转牌同花听牌面对底池下注}

\subsubsection*{牌局背景}

有效筹码 100BB。你在按钮位拿到：

\begin{center}
\texttt{As 7s}
\end{center}

翻牌为：

\begin{center}
\texttt{Ks 9s 2d}
\end{center}

翻牌你面对小注跟注。

转牌为：

\begin{center}
\texttt{4h}
\end{center}

大盲下注底池大小，行动轮到你。

\subsubsection*{新手常见想法}

很多新手会想：

\begin{itemize}
\item “我还是同花听牌，继续跟。”
\item “河牌来黑桃就赢。”
\item “翻牌都跟了，转牌不能弃。”
\end{itemize}

这正是常见的转牌追牌错误。

\subsubsection*{理性分析}

第一，数 outs。你仍然有 9 个同花 outs。

第二，估胜率。现在只剩河牌一张牌，用 2/4 法则：

\[
9 \times 2 = 18\%
\]

你完成同花的概率大约是 18\%。

第三，计算所需胜率。面对底池大小下注，你大约需要 33.3\% 胜率。

第四，比较。18\% 明显低于 33.3\%。从直接赔率看，跟注不划算。

第五，考虑隐含赔率。你是 A 高同花听牌，完成后很强。但如果河牌出现第三张黑桃，对手是否还会继续支付很多？如果不会，你的隐含赔率不足以弥补当前价格。

\subsubsection*{推荐行动}

多数情况下弃牌。

除非对手极有可能在河牌继续支付大量筹码，或者你有额外胜率和明确后续计划，否则面对转牌底池下注不应机械追听牌。

\subsubsection*{本手牌教训}

同样是同花听牌，翻牌和转牌价值完全不同。翻牌有两张牌可看，转牌只剩一张。转牌面对大注时，听牌继续条件更苛刻。

\subsection{牌例三：两头顺听牌面对小注}

\subsubsection*{牌局背景}

有效筹码 100BB。你在大盲拿到：

\begin{center}
\texttt{9c 8c}
\end{center}

按钮位 open raise 到 2.5BB，你跟注。

翻牌为：

\begin{center}
\texttt{7d 6s Kh}
\end{center}

你过牌，按钮下注 1/3 底池。

\subsubsection*{新手常见想法}

新手可能会想：

\begin{itemize}
\item “我还没成牌，但有顺子听牌。”
\item “下注不大，可以看一张。”
\item “来 5 或 T 就好了。”
\end{itemize}

这手牌可以用计算确认。

\subsubsection*{理性分析}

第一，数 outs。你需要 \texttt{5} 或 \texttt{T} 成顺。4 张 5 加 4 张 T，共 8 outs。

第二，估胜率。翻牌到河牌：

\[
8 \times 4 = 32\%
\]

第三，计算所需胜率。面对 1/3 底池下注，你大约需要 20\% 胜率。

第四，比较。32\% 高于 20\%。从直接赔率看，继续合理。

第五，考虑其他因素。你在大盲无位置，这是不利因素；但下注较小，听牌较清晰，继续可以接受。如果对手容易弃牌，也可以考虑加注半诈唬；如果对手不爱弃牌，跟注实现权益即可。

\subsubsection*{推荐行动}

可以跟注。根据对手类型，也可以考虑加注半诈唬。

\subsubsection*{本手牌教训}

两头顺听牌面对小注通常有足够价格继续。但无位置会降低实现权益的能力，因此不能只看 outs，还要考虑后续行动。

\subsection{牌例四：卡顺听牌面对大注}

\subsubsection*{牌局背景}

有效筹码 100BB。你在按钮位拿到：

\begin{center}
\texttt{Qc Jc}
\end{center}

你 open raise，大盲跟注。

翻牌为：

\begin{center}
\texttt{Ad Ks 4h}
\end{center}

大盲过牌，你下注 1/3 底池，大盲 check-raise 到接近底池大小加注。

\subsubsection*{新手常见想法}

新手可能会想：

\begin{itemize}
\item “我来 T 就是顺子。”
\item “我有卡顺，不能马上弃。”
\item “如果中了，可能赢大底池。”
\end{itemize}

但卡顺听牌只有 4 个 outs，面对大注必须非常谨慎。

\subsubsection*{理性分析}

第一，数 outs。你只有 \texttt{T} 能成顺，一共 4 张 T，因此是 4 outs。

第二，估胜率。翻牌到河牌：

\[
4 \times 4 = 16\%
\]

第三，计算所需胜率。面对接近底池大小的加注，你通常需要接近或超过 33\% 的胜率。

第四，比较。16\% 明显不够。

第五，考虑隐含赔率。虽然成顺后很强，但对手 check-raise 表示强范围，后续你未必中牌后能轻松拿到足够价值。此外，如果转牌没中，你很可能还要面对更大下注。

\subsubsection*{推荐行动}

弃牌。

\subsubsection*{本手牌教训}

卡顺听牌只有 4 outs，不能因为“中了很强”就面对大注继续。低 outs 听牌需要非常好的价格才值得追。

\subsection{牌例五：小对子追暗三条}

\subsubsection*{牌局背景}

有效筹码 150BB。你在按钮位拿到：

\begin{center}
\texttt{5c 5d}
\end{center}

前位玩家 open raise 到 3BB。你在按钮位行动。

\subsubsection*{新手常见想法}

新手可能会想：

\begin{itemize}
\item “我有对子，可以跟。”
\item “如果翻牌来 5，就很隐蔽。”
\item “不中就弃，也不会亏太多。”
\end{itemize}

这类手牌的价值主要来自隐含赔率。

\subsubsection*{理性分析}

第一，数 outs。如果你想从口袋对子变暗三条，牌堆中还剩 2 张 5，因此只有 2 outs。

第二，翻牌击中暗三条的概率并不高。粗略理解即可：大约八九次中一次。也就是说，大多数时候你会错过翻牌。

第三，为什么仍然可以跟？因为有效筹码深，你有位置，而且如果你中了暗三条，对手前位强范围可能有 \texttt{AA}、\texttt{KK}、\texttt{AK}、强顶对等愿意支付。这提供了隐含赔率。

第四，什么时候不能跟？如果有效筹码很浅，你即使中了也赢不到更多；如果后面还有激进玩家可能 squeeze；如果对手很紧且不中牌不会支付，那么跟注价值下降。

\subsubsection*{推荐行动}

在有效筹码较深、你有位置、后面再加注风险不高时，可以跟注。若筹码较浅或后面玩家激进，应更倾向弃牌。

\subsubsection*{本手牌教训}

小对子追暗三条不是靠直接赔率，而是靠隐含赔率。你不是因为经常中，而是因为一旦中牌可能赢大底池。

\subsection{牌例六：低同花听牌的反向隐含赔率}

\subsubsection*{牌局背景}

有效筹码 150BB。你在大盲拿到：

\begin{center}
\texttt{8s 6s}
\end{center}

按钮位 open raise，你跟注。

翻牌为：

\begin{center}
\texttt{Ks Qs 2d}
\end{center}

你过牌，按钮下注 2/3 底池。

\subsubsection*{新手常见想法}

新手可能会想：

\begin{itemize}
\item “我有同花听牌，9 outs。”
\item “如果来黑桃，我就是同花。”
\item “筹码深，中了可以赢很多。”
\end{itemize}

但这手牌有明显反向隐含赔率。

\subsubsection*{理性分析}

第一，表面 outs。你看似有 9 个黑桃 outs。

第二，outs 质量。你追的是 8 高同花。如果按钮有 \texttt{AsXs}、\texttt{JsTs}、\texttt{Ts9s} 等更高同花听牌，你完成同花后可能输大底池。

第三，位置问题。你在大盲无位置。即使转牌来黑桃，你也很难控制底池，并且对手如果强势下注，你会很尴尬。

第四，下注尺度。面对 2/3 底池下注，需要约 28.6\% 胜率。翻牌到河牌表面有约 36\%，但考虑 outs 不干净、无位置和反向隐含赔率，真实价值下降。

\subsubsection*{推荐行动}

可以谨慎跟注，但不能把它当作强听牌。面对更大下注、激进对手或后续难实现权益时，弃牌也完全合理。

\subsubsection*{本手牌教训}

低同花听牌的表面 outs 很好看，但并不等于真实价值很高。成牌后可能输给更高同花，是反向隐含赔率的典型来源。

\section{本章总结}

本章补充了德州扑克新手必须掌握的基础计算工具：outs、胜率和赔率。

Outs 是还能帮助你改善的补牌。常见同花听牌有 9 outs，两头顺听牌有 8 outs，卡顺听牌有 4 outs，一对变三条有 2 outs。数 outs 是判断听牌价值的第一步。

胜率是你最终赢下底池的大致概率。新手最实用的方法是 2/4 法则：

\begin{center}
\textbf{翻牌到河牌：outs $\times 4$。\\
转牌到河牌：outs $\times 2$。}
\end{center}

底池赔率告诉你，面对当前下注，你跟注需要赢多少比例才不亏。公式是：

\[
\text{所需胜率} =
\frac{\text{跟注金额}}{\text{跟注后总底池}}
\]

小注需要的胜率低，大注需要的胜率高。面对 1/3 底池下注约需 20\% 胜率，面对半池下注约需 25\%，面对底池下注约需 33.3\%。

真正的跟注判断，需要比较实际胜率和所需胜率：

\begin{center}
\textbf{实际胜率高于所需胜率，可以考虑继续。\\
实际胜率低于所需胜率，通常应该弃牌。}
\end{center}

但计算只是起点。实战中，你还必须考虑 outs 是否干净、是否有隐含赔率、是否存在反向隐含赔率、是否有位置、对手是否会支付、后续是否容易实现权益。

本章最重要的一句话是：

\begin{center}
\textbf{能不能跟，不看“有没有希望”，而看“这个希望值不值这个价格”。}
\end{center}

下一章“弃牌与跟注”将建立在本章基础之上，进一步讨论面对下注时如何估算对手价值牌和诈唬牌，哪些牌适合 bluff-catch，哪些看似强的牌应该弃，以及如何克服“不想被诈唬”的心理。

\section{课后训练}

\subsection*{训练一：Outs 速算}

找出 20 个翻牌或转牌听牌局面，逐一写下：

\begin{enumerate}
\item 我是什么类型的听牌？
\item 我有多少 outs？
\item 这些 outs 是否干净？
\item 如果不干净，原因是什么？
\end{enumerate}

训练目标是让你从“我有听牌”升级为“我有多少有效补牌”。

\subsection*{训练二：2/4 法则练习}

对以下常见 outs 进行胜率估算：

\begin{itemize}
\item 9 outs；
\item 8 outs；
\item 6 outs；
\item 4 outs；
\item 2 outs；
\item 15 outs。
\end{itemize}

分别计算：

\begin{enumerate}
\item 翻牌到河牌的大致概率；
\item 转牌到河牌的大致概率。
\end{enumerate}

训练目标是形成快速估算直觉。

\subsection*{训练三：底池赔率计算}

假设当前底池为 100，分别面对以下下注，计算你需要多少胜率：

\begin{itemize}
\item 对手下注 25；
\item 对手下注 33；
\item 对手下注 50；
\item 对手下注 75；
\item 对手下注 100；
\item 对手下注 150。
\end{itemize}

使用公式：

\[
\text{所需胜率} =
\frac{\text{跟注金额}}{\text{跟注后总底池}}
\]

训练目标不是追求小数点精确，而是熟悉下注越大、所需胜率越高。

\subsection*{训练四：听牌是否值得跟}

记录最近 10 手你用听牌跟注的手牌，逐手回答：

\begin{enumerate}
\item 我有多少 outs？
\item 我在翻牌还是转牌？
\item 我的估算胜率是多少？
\item 对手下注给我的所需胜率是多少？
\item 实际胜率是否高于所需胜率？
\item 我是否考虑了隐含赔率和反向隐含赔率？
\end{enumerate}

训练目标是减少无脑追听牌。

\subsection*{训练五：隐含赔率判断}

找出 5 手你追小对子、同花听牌或顺子听牌的手牌，分析：

\begin{itemize}
\item 有效筹码是否足够深？
\item 如果我成牌，对手是否会支付？
\item 我的成牌是否足够隐蔽？
\item 我是否有位置？
\item 如果这些条件不满足，我是否高估了隐含赔率？
\end{itemize}

训练目标是避免用“中了赢很多”作为模糊理由。

\subsection*{训练六：反向隐含赔率复盘}

找出 5 手你成牌后仍然输大底池的手牌，分析是否存在反向隐含赔率。

重点回答：

\begin{enumerate}
\item 我追的是不是非坚果牌？
\item 我完成后是否可能输给更强同花、更大顺子或更好踢脚？
\item 我是否把不干净 outs 算满了？
\item 我是否在成牌后过度自信？
\item 如果重来一次，我会不会在前面街更谨慎？
\end{enumerate}

训练目标是理解：不是所有成牌都值得追，也不是所有成牌后都适合打大底池。
